\begin{abstract}
Execution accuracy (EX), the standard Text-to-SQL metric, accepts a predicted query if it returns the gold result on one database; recent verifiers report that many accepted predictions are not equivalent to the gold query. We study what such counterexample databases prove, in three pre-registered stages on \V{nExc} accepted BIRD predictions of ten systems, thirteen further systems and Spider. Under one SQLite certifier, budget and set of admissibility rules, we compare five disprover families, including CaliberGraph, a join-graph-based disprover that certifies more pairs than the original SpotIt re-run under the same budget. Counts depend on the rules: almost no witness admissible under schema- and format-derived rules satisfies a data profile of the benchmark instance, although searching real sub-databases still certifies two thirds as many pairs. Whether a witness is evidence is harder to settle. An LLM annotator and two human annotators shown facts about the benchmark data judged most schema-tier witnesses to rest on implausible data, but in a third round three annotators agreed with each other barely above chance (Fleiss' $\kappa=\V{hvThreeFleissBin}$ on implausible vs.\ other), and their majority called only \V{hvThreeImpOne} of \V{hvThreeMajOne} such witnesses implausible. At the profile tier, our CaliberGraph-R certifies more predictions than random fuzzing on BIRD and on the new systems, but not on Spider. Counting every profile-tier witness as an error lowers EX by \V{ptsCountRMin}--\V{ptsCountRMax} points per system; weighting by attributed error rates gives \V{ptsErrRMin}--\V{ptsErrRMax} (LLM) or \V{ptsHumThreeRMin}--\V{ptsHumThreeRMax} points (human majority), both uncertain. We report every pre-registered outcome, including failures, and release all artifacts.
\end{abstract}
